\unnumberedchapter{KẾT LUẬN VÀ HƯỚNG PHÁT TRIỂN}

\section*{1. Kết quả đạt được}

Tiểu luận đã đi trọn vẹn qua bảy chủ đề của môn Phát triển hệ thống thông minh, từ các mô hình đơn lẻ trên bảng, ảnh, chuỗi, văn bản, đồ thị cho tới hệ đa tác tử. Mỗi chủ đề đều được kiểm chứng bằng thực nghiệm trên dữ liệu công khai về đái tháo đường, đối chiếu thêm trên các miền dữ liệu khác, cài đặt bằng nhiều cách (NumPy tự viết từ đầu, Keras, PyTorch) và đóng gói thành notebook cùng ứng dụng web cục bộ. \tabref{tab:summary} tóm tắt dữ liệu, mô hình tốt nhất và kết quả chính của cả bảy chương.

\begin{table}[H]
\centering
\caption{Tổng hợp kết quả thực nghiệm của bảy chương}
\label{tab:summary}
\small
\begin{tabularx}{\linewidth}{@{}cL{3.0cm}L{3.3cm}Y@{}}
\toprule
\textbf{Ch.} & \textbf{Dữ liệu} & \textbf{Mô hình tốt nhất} & \textbf{Kết quả chính (so với đối chứng)} \\
\midrule
1 & Pima (768) & ADALINE / hồi quy logistic & AUC 0,837; hệ chuyên gia 6 luật đạt 0,795; perceptron 0,749 và thất bại với XOR \\
2 & NHANES (20,6k) \& Diabetes (100k) & Rừng ngẫu nhiên; MLP 3 cách & NHANES AUC 0,823 (trần yếu tố nguy cơ); Diabetes 100k: MLP NumPy/Keras/PyTorch trùng khít AUC 0,977–0,978, F1 0,804–0,806 \\
3 & RetinaMNIST (1.600) \& MNIST/FMNIST & ResNet-18 chuyển giao; CNN 3 cách & RetinaMNIST: ResNet-18 QWK 0,700 so với 0,47–0,59 (CNN 3 cách học từ đầu) và 0,52 (MLP); trên MNIST CNN 3 cách đạt 98,8–99,7\% \\
4 & Shanghai CGM (112 người) \& tài chính & LSTM 2–5 tầng (3 cách cài đặt) & RMSE 30 phút 13,66–13,84 mg/dL (persistence 19,05; NumPy RNN 15,02); 95,5\% vùng Clarke A; giảm 27\% sai số trên CGM nhưng 0\% trên giá cổ phiếu/vàng \\
5 & KB ADA + MedQuAD + KG 96k nút & E5-large; GraphRAG / KG-RAG & Hit@1 KB tiếng Việt 0,967 (rerank), 0,933 (E5) so với 0,80 (BM25); RAG phủ 95\% dữ kiện, 0\% số liệu không căn cứ; KG-RAG tính chính xác thống kê trên 1.830 hồ sơ \\
6 & Pubmed-Diabetes (19.717) \& Cora/CiteSeer & GCN (ít nhãn), GraphSAGE (đủ nhãn) & PubMed 60 nhãn: GCN tự cài đặt đạt 0,789 so với 0,722 (MLP cùng 32.259 tham số); lợi ích đồ thị +6,7 điểm (PubMed) tới +24,4 điểm (Cora); GCN không kết nối tắt sụp từ 6 tầng \\
7 & Mèo–Chuột (8×8) \& MedMCQA (150) / GSM8K (40) & Q-learning đối kháng; Đa tác tử LLM có cổng đồng thuận & Bài 1: 2 mèo IQL bắt 99,7\%; mèo Q đơn tác tử sụp từ 100\% xuống 4\% trước chuột học, học đối kháng phục hồi lên 75\%. Bài 2: MAS đạt 72,5\% trên GSM8K (94,4\% khi CoT=PAL); trên MedMCQA cổng đồng thuận CoT=RAG đạt 70,1\% (Oracle 71,3\%), toàn tập 56,0\% do 23 lỗi phối hợp của Critic 1,5B \\
\bottomrule
\end{tabularx}
\end{table}

Từ các thực nghiệm xuyên suốt bảy chương có thể rút ra sáu bài học chung.

\emph{Thứ nhất, thông tin trong dữ liệu và giao thức đánh giá quyết định trần hiệu năng hơn là thuật toán hay framework.} Ở Chương 2,3,4, cùng một kiến trúc viết bằng NumPy, Keras hay PyTorch đều hội tụ về cùng một mức chất lượng (sai lệch đạo hàm chỉ $10^{-12}\sim 10^{-11}$), chỉ khác nhau tốc độ thực thi. Ở Chương 4, cùng mạng LSTM giảm 27–35\% sai số trên chuỗi glucose và doanh thu nhưng không thắng nổi dự báo ngây thơ trên giá cổ phiếu và vàng.

\emph{Thứ hai, thiên kiến quy nạp phù hợp với cấu trúc dữ liệu tạo nên khác biệt lớn nhất.} Tích chập (kết nối cục bộ, chia sẻ trọng số) cho ảnh, cổng nhớ cộng tính cho chuỗi, phép nhân ma trận kề chuẩn hóa $\hat A$ (không tốn thêm tham số nào) cho đồ thị: mỗi kiến trúc thắng mô hình ``phẳng'' ở đúng loại dữ liệu của nó với ít tham số hơn.

\emph{Thứ ba, tri thức có sẵn rất có giá trị khi dữ liệu nhỏ.} Học chuyển giao từ ImageNet nâng QWK thêm khoảng 0,15. Một bộ 6 luật y khoa đánh bại perceptron. Đồ thị trích dẫn bù cho việc chỉ có 20 nhãn mỗi lớp.

\emph{Thứ tư, mô hình đơn giản là mốc đối chứng không được bỏ qua.} Persistence, hồi quy tuyến tính, BM25 và trả lời trực tiếp (\texttt{Direct}) trên câu hỏi sự kiện đều mạnh hơn dự kiến; chỉ những cải thiện vượt được chúng mới thực sự có ý nghĩa.

\emph{Thứ năm, trong hệ đa tác tử, môi trường không dừng và cơ chế phối hợp là hai nút thắt quyết định.} Với học tăng cường, chính sách học trước đối thủ cố định sẽ quá khớp với đối thủ đó (tỷ lệ bắt sụt từ 100\% xuống 4\%) nếu không được huấn luyện đối kháng. Với tác tử LLM, kiểm chứng chéo hai cơ chế độc lập (cổng đồng thuận \texttt{CoT = RAG} hoặc \texttt{CoT = Python}) nâng mạnh độ tin cậy (70,1\% và 94,4%), nhưng nếu dùng một LLM nhỏ thiếu công cụ tất định làm trọng tài (\texttt{Critic}) khi bất đồng thì lỗi phối hợp sẽ làm hao hụt phần lớn lợi thế của tập hợp tác tử.

\emph{Thứ sáu, trong y tế, hệ thống tốt là hệ thống biết giới hạn của mình và kiểm chứng được.} Ngưỡng xét nghiệm chuẩn (ADA) vẫn là căn cứ chẩn đoán; mô hình học được chỉ bổ sung ước lượng nguy cơ. Với RAG và đa tác tử, mọi câu trả lời cần kèm trích dẫn, vết suy luận minh bạch và biết từ chối hoặc báo bất đồng khi thiếu căn cứ.

\section*{2. Hạn chế}

\begin{itemize}
  \item Dữ liệu huấn luyện chủ yếu đến từ Hoa Kỳ (NHANES, Pima), Trung Quốc (Shanghai) và Ấn Độ (MedMCQA). Chưa có dữ liệu lâm sàng thật của người Việt Nam; nhóm người gốc Á của NHANES chỉ là phép xấp xỉ.
  \item RetinaMNIST có độ phân giải thấp (128×128) và chỉ 1.600 ảnh. Mô hình CNN chưa đủ tin cậy để dùng lâm sàng; Grad-CAM cho thấy mô hình đôi khi chú ý vào vùng lóa và viền ảnh.
  \item Dự báo glucose chưa dùng thông tin bữa ăn và liều tiêm insulin, nên sai số tăng nhanh sau 30 phút.
  \item Ở Chương 7, giới hạn 4 GB VRAM buộc hệ đa tác tử phải dùng mô hình cục bộ kích thước nhỏ (\texttt{Qwen2.5-1.5B-Instruct}) và kho văn bản phổ thông MedQuAD-NIDDK, khiến tác tử \texttt{Critic} chưa đủ thẩm quyền y khoa để phân xử các câu hỏi mức sau đại học khi hai solver bất đồng (gây ra 23 lỗi phối hợp).
\end{itemize}

\section*{3. Hướng phát triển}

\begin{itemize}
  \item Thu thập dữ liệu sàng lọc của người Việt (phối hợp cơ sở y tế) để hiệu chỉnh lại mô hình và ngưỡng quyết định.
  \item Thay RetinaMNIST bằng ảnh đáy mắt độ phân giải cao (APTOS, EyePACS), áp dụng mô hình nền tảng thị giác cho ảnh y khoa.
  \item Dự báo glucose đa biến (bữa ăn, insulin, vận động) bằng Transformer cho chuỗi thời gian; dùng hàm mất mát bất đối xứng ưu tiên vùng hạ đường huyết.
  \item Ở hệ đa tác tử y khoa (Chương 7), thay bộ phân xử thuần ngôn ngữ của LLM 1,5B bằng bộ kiểm chứng lai kết hợp Đồ thị tri thức lâm sàng (Chương 5) và công cụ kiểm tra ngưỡng xét nghiệm bằng Python, từ đó thu hẹp khoảng cách từ 56,0\% tiến tới cận trên Oracle 71,3\%.
  \item Chuẩn hóa giao tiếp công cụ theo Model Context Protocol \cite{mcp2024}; thêm đăng nhập, phân quyền và nhật ký kiểm toán để đáp ứng yêu cầu bảo vệ dữ liệu sức khỏe.
\end{itemize}
